\begin{figure}[H]
\centering
\begin{tikzpicture}[>=Stealth,font=\small]
\node at (0,1.6) {$\BaseCarrier$};
\draw (0,0) ellipse (1.15 and 0.4);
\fill (-0.55,0.04) circle (1.1pt);
\fill (0.5,-0.12) circle (1.1pt);
\fill (0.05,0.08) circle (1.5pt);
\node[above,inner sep=3pt] at (0.05,0.08) {$\BasePoint$};
\draw[->] (1.5,0) -- node[above]
 {$\times\EquilateralCarrier_3$} (3.1,0);
\node at (4.65,1.6) {$\BaseCarrier\times\EquilateralCarrier_3$};
\fill[black!8] (4.5,-1.15) rectangle (4.9,1.3);
\foreach \level/\index in {0.95/0,0/1,-0.95/2}
{
 \draw (4.65,\level) ellipse (1.15 and 0.4);
 \fill (4.1,{\level+0.04}) circle (1.1pt);
 \fill (5.15,{\level-0.12}) circle (1.1pt);
 \fill (4.7,{\level+0.08}) circle (1.5pt);
 \node[right] at (5.95,\level)
  {$\BaseCarrier\times\{\index\}$};
}
\node[align=center] at (4.7,-1.95)
 {$\{\BasePoint\}\times\EquilateralCarrier_3$\\
  pairwise distance $\ProductScale$};
\end{tikzpicture}
\par\medskip
\begin{tikzpicture}[
 >=Stealth,
 font=\small,
 role/.style={draw,align=center,text width=3.4cm,
              minimum height=2.1cm,inner sep=5pt,anchor=north}
]
\node at (4,0.9)
 {$\MetricSpace{\ClosedDomain}{\SubspaceMetric}
   =\ProductApproximation_{\EquilateralSize_{\FamilyIndex}}
      \MetricSpace{\BaseCarrier}{\MetricSymbol},
   \qquad
   \ProductScale=\CoverScale\MetricSpace{\BaseCarrier}{\MetricSymbol}$};
\node (formula) at (4,0)
 {$\MetricSpace{\ClosedDomain}{\SubspaceMetric}
   =\MetricSpace{\BaseCarrier\times\EquilateralCarrier_{\EquilateralSize_{\FamilyIndex}}}
                {\MetricSymbol\vee\EquilateralMetric_{\EquilateralSize_{\FamilyIndex},\ProductScale}}$};
\node[role] (approximation) at (0,-1.7)
 {small scale $\ProductScale$\\
  $\begin{gathered}
    \GHDistance(\MetricSpace{\ClosedDomain}{\SubspaceMetric},
                 \MetricSpace{\BaseCarrier}{\MetricSymbol})\\
    \leq\ProductScale/2
   \end{gathered}$\\
  approximation};
\node[role] (separation) at (4,-1.7)
 {$\EquilateralSize_{\FamilyIndex}>\PackingBound_{\FamilyIndex}$\\
  recursive choice\\
  disjoint images};
\node[role] (packing) at (8,-1.7)
 {$\EquilateralSize_{\FamilyIndex}\to\infty$\\
  positive scale bound\\
  no accumulation};
\draw[->,shorten <=2pt]
 ([xshift=-1.8cm]formula.south) .. controls +(0,-0.55) and +(0,0.55) .. (approximation.north);
\draw[->,shorten <=2pt]
 (formula.south) -- (separation.north);
\draw[->,shorten <=2pt]
 ([xshift=1.8cm]formula.south) .. controls +(0,-0.55) and +(0,0.55) .. (packing.north);
\end{tikzpicture}
\caption{The upper diagram shows three copies of
$\BaseCarrier$
in a product with an equilateral space.
The marked fiber has
$\EquilateralSize$
points at pairwise distance
$\ProductScale$
for a general factor
$\EquilateralCarrier_\EquilateralSize$.
The scales and cardinalities of the equilateral factors control the
 discrete approximation.
The bound
$\PackingBound_{\FamilyIndex}$
applies to subsets separated by
$\ProductScale_{\FamilyIndex}/2$
in the previously constructed images,
where
$\ProductScale_{\FamilyIndex}$
is the minimum input scale in \eqref{eq:minimum-product-scale}.
The recursive choice \eqref{eq:equilateral-cardinality-choice}
separates the images.
Along a convergent sequence of approximating spaces,
the scales have a positive lower bound,
so the divergence of the cardinalities contradicts
Corollary~\ref{cor:uniform-separated-sets}.}
\label{fig:discrete-approximation}
\end{figure}
